## TEMP holding note — does RQ2a's own residual-reduction value cluster spatially? (2026-08-16)

**`[SUPERSEDED 2026-08-18]`**: k=8 nearest-neighbour method replaced project-wide by a
semivariogram-informed distance-band method — same change as `TEMP_rq2_attribution_results_2026-08-11.tex`
and `TEMP_rq3_gnnwr_results_2026-08-11.tex`. **The conclusion below does not just shrink in
magnitude, it reverses.** Corrected numbers and conclusion follow; original k=8 table kept inline
for provenance, not for citing.

**Context**: during figure planning, a spatial "help map" for RQ2a (where does environmental
conditioning help or hurt) was explicitly NOT proposed, because no supporting spatial-clustering
statistic existed yet for the reduction values themselves — only for raw residuals, computed
separately for RQ2b/RQ3. This closes that gap.

**Method (re-run 2026-08-18)**: `models/spatial_attribution/rq2a_reduction_morans_i.py`. No new
fitting — composes three already-existing pieces: `compute_residual_reduction()`
(`models/spatial_attribution/rq2_residual_reduction.py`, the exact function already behind RQ2a's
own quartile tables) for the per-row reduction value (`|CR error| - |model error|`);
`load_plot_coordinates()` (`models/common/geo.py`) for (x, y); `compute_residual_morans_i()`
(`models/growth_curve_attribution/residual_spatial_autocorrelation_check.py`, now
semivariogram-informed distance-band weights, 999 permutations, seed=42, max_distance=5000m,
4survey re-checked at 10000m since neither model resolved cleanly at 5000m — see below) for the
test itself. DNN and XGBoost (both env-conditioned, Set3, seed 42/fixed config), both cohorts,
pooled across all 5 `spatial_block_kfold` test folds.

### Results

| Cohort | Model | n | Moran's I @5000m | status @5000m | Moran's I @10000m | status @10000m |
|---|---|---:|---:|---|---:|---|
| 4survey | DNN | 232,292 | +0.0380 | resolved (range 4802m, 198m from ceiling) | -0.0038 | exceeds_window |
| 4survey | XGBoost | 232,292 | +0.0401 | exceeds_window (5000m) | -0.0045 | exceeds_window |
| 6survey | DNN | 82,614 | NaN | no_structure | not re-checked (no_structure is a null, not a narrow-window artifact) | — |
| 6survey | XGBoost | 82,614 | NaN | no_structure | not re-checked | — |

**The reduction value shows no reliable spatial clustering, for either model, on either cohort —
the opposite of the original k=8 finding.** 6survey: semivariogram flat even at the 5000m window on
both models, a genuine null. 4survey: the 5000m-window DNN value (+0.038) looked real but its range
(4802m) sat suspiciously close to the search ceiling; widening to 10000m, neither model resolves
even then, and both flip sign toward zero (DNN +0.038→-0.004, XGBoost +0.040→-0.005) — the
signature of no genuine spatial structure, not of structure at a wider scale. This remains an
`exceeds_window` (unresolved) result at 10000m, not a formally proven null, but combined with
6survey's solid null it does not support real spatial structure.

**What this does and doesn't establish**: does not confirm that where conditioning helps is
spatially clustered — on the evidence available, it is not. This is a different, separate claim
from the already-established CR-residual-quartile mechanism (RQ2a item 1: conditioning helps where
CR fits worst, hurts where it fits well), which is untouched by this correction — that finding
does not depend on spatial clustering of the reduction value. No causal claim is made either way.

**Implication for the plotting plan — reversed**: the planned spatial "help map" for RQ2a is **not
supported** by this corrected check. Do not build it on the basis of this analysis; the original
2026-08-16 justification ("this now supports a genuine, evidence-based spatial help map") does not
hold under the corrected method.

**Original k=8 table (2026-08-16), for provenance only, not for citing**: 4survey DNN +0.6424, XGB
+0.6281; 6survey DNN +0.5698, XGB +0.5499 (all p=0.001, same n as above).
